\documentclass[12pt]{article}

\input{../../_assets/preambulo.tex}


\begin{document}

    % 1. Foto de fondo
    % 2. Título
    % 3. Encabezado Izquierdo
    % 4. Color de fondo
    % 5. Coord x del titulo
    % 6. Coord y del titulo
    % 7. Fecha

    
    \input{../../_assets/portada}
    \portadaExamen{ffccA4.jpg}{Ecuaciones\\Diferenciales II\\Examen XV}{Ecuaciones Diferenciales II. Examen XV}{MidnightBlue}{-8}{28}{2026}{}

    \begin{description}
        \item[Asignatura] Ecuaciones Diferenciales II.
        \item[Curso Académico] 2022/2023.
        \item[Grado] Doble Grado en Ingeniería Informática y Matemáticas.
        \item[Grupo] Único.
        \item[Profesor] Rafael Ortega Ríos.
        \item[Descripción] Convocatoria Ordinaria.
        \item[Fecha] 16 de junio de 2023.
        % \item[Duración] ---.
    \end{description}
    \newpage


    % ------------------------------------
    
    \begin{ejercicio}
        Para cada $n \geq 1$ la función $f_n : \left[-1, 1\right] \longrightarrow \bb{R}$ es continua y cumple $f_n\left(\nicefrac{1}{n}\right) = f_n\left(\nicefrac{-1}{n}\right) = 1$, $f_n(0) = 0$. Demuestra que la sucesión $\{f_n\}_{n \geq 1}$ no converge uniformemente.
    \end{ejercicio}

    \begin{ejercicio}
        Encuentra todas las funciones continuas $x : \left[0, 1\right] \longrightarrow \bb{R}$ que cumplen
        \[
        x(t) = \int_{0}^{\nicefrac{t}{2}} x(s) \, ds, \quad \forall t \in \left[0, 1\right].
        \]
    \end{ejercicio}

    \begin{ejercicio}
        Dados $t_0, x_0 \in \bb{R}$, encuentra el intervalo maximal de la solución del problema
        \[
        \dot{x} = x^3, \quad x(t_0) = x_0.
        \]
    \end{ejercicio}

    \begin{ejercicio}
        Se considera la ecuación integral
        \[
        x(t) = \operatorname{sen}\left(2 + \int_{0}^{t} x(s) \, ds\right), \quad t \in \bb{R}.
        \]
        Se supone que las posibles soluciones son funciones continuas $x : \bb{R} \longrightarrow \bb{R}$. Encuentra un problema de valores iniciales que sea equivalente.
        \begin{observacion}
            La incógnita será una primitiva de $x(t)$.
        \end{observacion}
    \end{ejercicio}

    \begin{ejercicio}
        La solución del problema de Cauchy
        \[
        \ddot{\theta} + 4 \operatorname{sen} \theta = 0, \quad \theta(0) = \theta_0, \quad \dot{\theta}(0) = \omega_0
        \]
        se denota por $\theta(t; \theta_0, \omega_0)$. Calcula, si existe, $\del{\theta}{\omega_0}(t; \pi, 0)$.
    \end{ejercicio}

    \begin{ejercicio}~
        \begin{enumerate}[label=\roman*)]
            \item Dado un parámetro $\veps > 0$, $x(t, \veps)$ es la solución del problema
            \[
            \dot{x} = \operatorname{sen}(\veps x) + e^t, \quad x(0) = \pi.
            \]
            Calcula, si existe, $\lim\limits_{\veps \to 0^+} x(2\pi, \veps)$.
            \item Dado un parámetro $\veps > 0$, $x(t, \veps)$ es la solución del problema
            \[
            \veps \dot{x} = e^x \operatorname{sen} t, \quad x(0) = \pi.
            \]
            Calcula, si existe, $\lim\limits_{\veps \to 0^+} x(2\pi, \veps)$.
        \end{enumerate}
    \end{ejercicio}

    % \newpage
    % \setcounter{ejercicio}{0} % Reiniciar contador de ejercicios
    % \noindent
    % \textbf{Solución.}

\end{document}